\documentclass[../cours_3e.tex]{subfiles}
\begin{document}
\setcounter{chapter}{7}
\chapter{Comprendre les fonctions}
\label{ch:fonctions}
\begin{objectif}
Relier formule, tableau, programme et graphique ; calculer une image ; chercher des antécédents ; modéliser une grandeur. Prérequis : repérage et équations.
\end{objectif}
\section{Une entrée, une sortie}
\begin{definition}
Une \textbf{fonction} associe à chaque nombre autorisé en entrée un unique nombre en sortie. On écrit $f:x\mapsto f(x)$. Le nombre $f(x)$ est l'\textbf{image} de $x$. Si $f(a)=b$, $a$ est un \textbf{antécédent} de $b$.
\end{definition}
La fonction est le procédé entier ; $f(x)$ est le résultat pour une entrée donnée. Le $x$ est une variable. L'écriture $f(x)$ ne signifie pas $f\times x$.
\begin{exemple}
Si $f(x)=x^2-1$, alors $f(3)=8$ et $f(-3)=8$. Le nombre $3$ possède une seule image, mais $8$ possède deux antécédents : $3$ et $-3$.
\end{exemple}
Un nombre peut avoir aucun, un ou plusieurs antécédents. Ici $-2$ n'en a aucun, car $x^2-1\geq-1$.
\section{Le domaine des valeurs autorisées}
Le contexte ou la formule limite les entrées. Pour $g(x)=1/(x-2)$, $x=2$ est interdit. Pour une longueur, les valeurs négatives sont exclues. Pour un nombre de billets, seules des valeurs entières positives ou nulles conviennent.
\begin{attention}
Avant de chercher un antécédent, préciser l'intervalle ou l'ensemble étudié. Sur $x\geq0$, la fonction $x^2-1$ n'associe $8$ qu'à $3$ ; sur tous les réels, elle l'associe aussi à $-3$.
\end{attention}
\section{Quatre façons de décrire le même objet}
Pour $f(x)=x^2-1$ :
\begin{itemize}
\item formule : $f(x)=x^2-1$ ;
\item programme : choisir $x$, l'élever au carré, soustraire $1$ ;
\item tableau : calculer des valeurs choisies ;
\item graphique : placer les points de coordonnées $(x;f(x))$.
\end{itemize}
\[
\begin{array}{c|rrrrrrr}
x&-3&-2&-1&0&1&2&3\\\hline
f(x)&8&3&0&-1&0&3&8
\end{array}
\]
\begin{center}
\begin{tikzpicture}
\begin{axis}[width=10cm,height=7cm,axis lines=middle,xmin=-3.5,xmax=3.5,ymin=-2,ymax=9,grid=both,xtick={-3,-2,...,3},ytick={-1,0,1,3,5,8},xlabel={$x$},ylabel={$f(x)$}]
\addplot[accent,thick,domain=-3.15:3.15,samples=100]{x^2-1};
\addplot[only marks,mark=*,accent]coordinates{(-3,8)(-2,3)(-1,0)(0,-1)(1,0)(2,3)(3,8)};
\addplot[dashed,domain=-3.3:3.3]{3};
\end{axis}
\end{tikzpicture}
\end{center}
Un tableau ne donne généralement qu'une partie des valeurs. Plusieurs fonctions peuvent partager ces valeurs et différer entre elles. Relier arbitrairement les points par des segments ajoute donc une hypothèse injustifiée.
\section{Lire une image et un antécédent}
\begin{methode}[Lecture graphique]
Pour l'image de $a$, partir de $a$ sur l'axe horizontal, rejoindre verticalement la courbe, puis lire l'ordonnée. Pour les antécédents de $b$, tracer mentalement l'horizontale d'ordonnée $b$ et lire les abscisses de tous ses points d'intersection avec la courbe.
\end{methode}
Sur la figure, l'image de $2$ est $3$. Les antécédents de $3$ sont $-2$ et $2$. Si le graphique est une représentation approchée ou si la lecture ne tombe pas sur un point explicitement connu, on donne un résultat approché.
\begin{exemple}[À partir d'un tableau]
Dans le tableau précédent, $f(-1)=0$. Les valeurs $-1$ et $1$ sont des antécédents de $0$. Le tableau seul ne prouverait pas qu'il n'en existe aucun autre ; la formule $x^2-1=0$ le prouve.
\end{exemple}
\section{Calculer et résoudre}
\begin{methode}
Pour une image, remplacer $x$ par la valeur donnée. Pour un antécédent de $b$, résoudre $f(x)=b$ puis vérifier les contraintes.
\end{methode}
Si $h(x)=4x-7$, l'image de $-2$ est $-15$. L'antécédent de $9$ vérifie $4x-7=9$, donc $x=4$.
Pour $u(x)=(x-3)(2x+1)$, les antécédents de $0$ sont $3$ et $-1/2$ par la règle du produit nul.
\section{Modéliser un phénomène}
\begin{exemple}[Aire variable]
Un rectangle a un périmètre de $24\unite{cm}$. Si un côté mesure $x\unite{cm}$, l'autre mesure $12-x$. Son aire est
\[A(x)=x(12-x),\qquad 0<x<12.\]
$A(4)=32\unite{cm^2}$. Les valeurs $x=0$ et $x=12$ donnent des figures aplaties, exclues si l'on demande un vrai rectangle.
\end{exemple}
La fonction exprime une variation continue : le côté peut prendre toutes les longueurs réelles comprises entre $0$ et $12$. Pour des objets comptés, on ne retient au contraire que des entrées entières.
\begin{exemple}[Une même aire, plusieurs rectangles]
$A(4)=A(8)=32$ : échanger longueur et largeur ne change pas l'aire. Cette égalité s'explique aussi par $x(12-x)=(12-x)x$.
\end{exemple}
\utilite{Une fonction rassemble dans un même modèle une table de mesures, une formule et un graphique. On peut alors prévoir une valeur, chercher un seuil ou comparer deux procédés.}
\begin{retenir}
Image : partir de l'entrée. Antécédent : retrouver une entrée possible. Un point $(a;b)$ appartient à la courbe exactement lorsque $b=f(a)$. Toujours distinguer une lecture approchée d'un calcul exact.
\end{retenir}
\end{document}
